% ---------------------------------------------------------------------
% STATUT : RÉÉCRIT (juillet 2026). La thèse passe de « W n'est pas
%          identifiable » à une frontière en TROIS NIVEAUX : structure
%          corroborée, direction par paire non établie, amplitude jamais
%          identifiable. Sources : scripts 40, 53, 55, 56.
% LABELS À PRÉSERVER : ch:identifiabilite (cité par 8 chapitres),
%          fig:reversibilite (cité par 12b).
% ---------------------------------------------------------------------
\chapter{L'identifiabilité de la contagion dirigée}
\label{ch:identifiabilite}

\begin{cle}
\textbf{Thèse du chapitre.} La \emph{direction} de la contagion entre les cinq piliers DORA
n'est pas identifiable sur les données publiques agrégées disponibles, quand sa co-occurrence
l'est. C'est un résultat d'identification empirique : les trois sources échouent pour des raisons
propres à leur contenu, et la théorie des processus réversibles en donne la raison de principe
pour les statistiques invariantes par renversement du temps. Il ne vaut pas pour toute donnée, la
loi complète des piliers touchés par incident portant de l'information sur la direction
(proposition~\ref{prop:nonid}). Mais l'affaire ne s'arrête pas là : un corpus de rapports
d'incidents officiels corrobore la \emph{structure} de cette direction, sans lever pour autant
le doute sur chacun de ses coefficients. On aboutit donc à une frontière en trois
niveaux, que ce chapitre établit puis assume : la structure est revendiquée, la direction paire
par paire ne l'est pas, l'amplitude ne le sera jamais par cette voie. Le capital est en
conséquence lu en bornes, et la donnée manquante devient une spécification de reporting
chiffrée et hiérarchisée.
\end{cle}

\section{Ce qu'identifier $W$ exigerait}

Le coefficient $W_{jk}$ mesure la propension d'une défaillance du pilier $j$ à
\emph{entraîner}, dans la foulée, une défaillance du pilier $k$ (convention
du chapitre~\ref{chap:cascade} : la ligne émet, la colonne reçoit). L'identifier sur données
demande trois choses simultanément : (i) observer des co-occurrences de défaillances de
piliers au sein d'une même entité ; (ii) disposer d'une antériorité temporelle fiable
(qui précède qui) pour distinguer $W_{jk}$ de $W_{kj}$ ; (iii) une taxonomie par pilier
(domaine de contrôle), et non par vecteur d'attaque. Aucune des trois sources agrégées examinées
ne réunit les trois, et elles échouent pour des raisons \emph{différentes} : la première n'a pas
la taxonomie, la deuxième n'a pas le pas de temps, la troisième a les deux dans son schéma mais
ne les renseigne pas.

\section{Trois échecs, pour trois raisons distinctes}

\paragraph{Source 1 : Data Breach Chronology (notifications d'États américains).}
La base est volumineuse ($74\,783$ incidents, $25\,009$ organisations) mais inexploitable
pour $W$, par manque de taxonomie et de volume informatif. La catégorie
non déterminée (\texttt{UNKN}) pèse 52,7\,\% des incidents et n'est attribuable à
aucun pilier. La date de \emph{survenance} n'est exploitable que dans 49,6\,\% des
cas, et son caractère manquant n'est pas aléatoire : il varie par type d'incident, si bien
qu'ordonner par survenance élimine des catégories entières du panel. Ordonner par
\emph{déclaration}, l'autre date disponible, ne mesure pas la contagion mais le
délai de détection. Résultat : une fois restreint aux types informatifs et aux
trimestres consécutifs, il ne reste que 71 transitions $k(t\!-\!1)\to j(t)$ (soit
$3{,}6$ observations par cellule hors-diagonale), et 34 dans le seul secteur
financier. La direction n'y est pas estimable, faute de matière.

\paragraph{Source 2 : SAS OpRisk Global Data (pertes opérationnelles, Bâle).}
Ici le volume est au rendez-vous : 4\,329 transitions annuelles entre les
7 catégories de risque, soit environ 103 observations par cellule. Et pourtant la
direction reste inidentifiable, pour la raison \emph{opposée} : le pas de temps est
l'année, trop grossier pour porter une antériorité. Un test placebo le montre sans
appel. Soit $M$ la matrice des transitions observées, $M_{jk}$ comptant les passages
$k(t-1)\to j(t)$. On mesure la direction par la statistique d'asymétrie
\begin{equation}
  T(M) \;=\; \tfrac{1}{2}\,\lVert M - M^{\top}\rVert_1
        \;=\; \tfrac{1}{2}\sum_{j\neq k}\bigl|M_{jk}-M_{kj}\bigr| ,
  \label{eq:asymetrie}
\end{equation}
nulle si et seulement si $M$ est symétrique. Sous l'hypothèse nulle $H_0$ d'absence de
direction, permuter les années \emph{à l'intérieur de chaque firme} détruit toute antériorité
tout en conservant la co-occurrence : si $\Pi$ désigne une telle permutation et $M^{\Pi}$ la
matrice recalculée, la loi de $T(M^{\Pi})$ est celle de $T(M)$ sous $H_0$. En notant $\mu_0$ et
$\sigma_0$ ses moyenne et écart-type empiriques sur les tirages de $\Pi$, la statistique
centrée réduite
\begin{equation}
  z \;=\; \frac{T(M)-\mu_0}{\sigma_0}
  \label{eq:placebo-z}
\end{equation}
vaut ici
\[
T(M) = 119 \qquad\text{contre}\qquad \mu_0 \pm \sigma_0 = 128 \pm 27
\qquad\Longrightarrow\qquad z = -0{,}33 .
\]
L'asymétrie observée est \emph{sous} le bruit du placebo : aucun signal directionnel. Le
test binomial par paire confirme : sur 18 paires suffisamment peuplées, aucune
n'est significative à 5\,\%, ni avant ni après correction de Bonferroni. Le pas annuel a
lessivé la direction.

\begin{cle}
\textbf{Précaution de méthode (à afficher).} Le test intuitif du \og temps inversé \fg{}
(comparer $M$ à la matrice obtenue en renversant la chronologie) est dégénéré sur
un comptage : par construction $M_{\text{inv}} = M^\top$, donc la corrélation des asymétries
vaut $-1$ quel que soit le contenu. Il ne prouve rien. Seul le placebo par
permutation est valide sur des comptages.
\end{cle}

\paragraph{Source 3 : la VERIS Community Database, ou le schéma sans le remplissage.}
Les deux premiers échecs tiennent à ce que la source ne \emph{prévoit} pas ce que $W$ exige.
Restait l'hypothèse la plus favorable : un standard de description d'incidents qui, lui,
prévoit exactement les bons champs. C'est le cas du schéma VERIS, référence de la
profession et support du \emph{Data Breach Investigations Report}, dont la base publique
compte $10\,591$ incidents ($980$ en secteur financier, NAICS~52). Son schéma porte
littéralement les deux objets manquants : un champ \texttt{control\_failure}, qui nomme le
domaine de contrôle défaillant, et une chronologie en quatre étapes
(compromission, exfiltration, découverte, confinement), qui porte l'antériorité.

Le verdict est sans appel, et il porte sur le \emph{remplissage} et non sur la conception.
Le champ \texttt{control\_failure} est renseigné dans $13$ incidents sur $10\,591$
($0{,}12\,\%$), dont un seul en secteur financier. La chronologie n'est datée sur au
moins deux étapes que dans $126$ incidents ($1{,}19\,\%$), et sur au moins trois dans $10$.
Surtout, l'intersection des deux conditions, qui est ce qu'identifier $W$ exigerait
simultanément, ne compte que $1$ incident sur toute la base, et aucun en
secteur financier.

\begin{cle}
\textbf{Pourquoi ce troisième échec est le plus utile des trois.} Les deux premiers pouvaient
se lire comme un défaut de conception des sources : personne n'avait prévu de collecter la
direction. Celui-ci interdit cette lecture. La spécification \emph{existe}, elle est publiée,
elle est adoptée par la profession, et elle n'est pas remplie. Le problème n'est donc
pas de savoir \emph{quoi} collecter, question à laquelle le schéma VERIS répond déjà,
mais d'en rendre le remplissage obligatoire. C'est exactement ce qu'un registre
d'incidents sous DORA peut imposer et qu'une base déclarative volontaire ne peut pas, et c'est
ce qui donne sa portée à la spécification de reporting du chapitre~\ref{ch:livrable} : elle
ne demande pas d'inventer des champs, elle demande d'en rendre deux obligatoires.
\end{cle}

% SOURCES-SCRIPTS: 05 57

% SOURCES-SCRIPTS: 28 40 64 110

% SOURCES-SCRIPTS: 40 64

\section{Pourquoi la direction échappe à la co-occurrence : la lecture par la réversibilité}

Le placebo établit un fait ; la théorie des processus réversibles en donne la raison de principe,
sous une hypothèse qu'il faut écrire : que les sauts entre piliers se représentent par une chaîne
de Markov stationnaire. Toute matrice se décompose de façon unique en $W=S+A$, part
symétrique $S=(W+W^\top)/2$ et antisymétrique $A=(W-W^\top)/2$, et ces deux parts sont
\emph{orthogonales} au sens de Frobenius. Sur la chaîne de Markov des sauts inter-piliers construite à partir
du jugement dirigé (de loi stationnaire $\pi$), $S$ est la partie réversible et $A$ le
\emph{courant} de probabilité $J_{jk}=\pi_j P_{jk}-\pi_k P_{kj}$, antisymétrique et nul si et
seulement si l'équilibre détaillé est satisfait. Le renversement du temps laisse $S$ invariante
et retourne le signe de $A$.

La décomposition $W=S+A$ est ce qui porte les bornes du chapitre~\ref{chap:partielle} et elle se
retient telle quelle. Les trois identités qui suivent ne portent, elles, aucune valeur publiée :
elles établissent pourquoi une statistique invariante par renversement du temps ne \emph{pouvait}
pas révéler la direction, donc pourquoi l'échec des trois sources tient à la nature de
l'observable et non à un défaut de collecte. Elles sont empruntées à la théorie des chaînes de
Markov et non à l'outillage actuariel, et les sauter ne coûte rien à la suite du mémoire.

Formellement, en notant $P$ le noyau de transition et $\pi$ sa loi stationnaire, on décompose
$P = S_\pi + A_\pi$ où $S_\pi=\tfrac12(P+\tilde P)$ et $A_\pi=\tfrac12(P-\tilde P)$, avec
$\tilde P_{jk} = \pi_k P_{kj}/\pi_j$ le noyau \emph{adjoint} (la chaîne renversée dans le temps).
Trois identités s'ensuivent, vérifiées numériquement (%
% SOURCES-SCRIPTS: 40
figure~\ref{fig:reversibilite}).

\emph{(i) Le courant.} La quantité
\begin{equation}
  J_{jk} \;=\; \pi_j P_{jk} - \pi_k P_{kj}
  \label{eq:courant}
\end{equation}
est antisymétrique, et $J\equiv 0$ caractérise la réversibilité (équilibre détaillé). Elle est
ici non nulle : le modèle porte bien une direction.

\emph{(ii) La fluctuation ne voit que la partie réversible.} Pour toute fonction test
$f:\{1,\dots,5\}\to\R$, la forme de Dirichlet, qui est la variation quadratique de la martingale
associée, vaut
\begin{equation}
  \mathcal{E}(f,f) \;=\; \tfrac12\sum_{j,k}\pi_j P_{jk}\,(f_k-f_j)^2
                  \;=\; \tfrac12\sum_{j,k}\pi_j (S_\pi)_{jk}\,(f_k-f_j)^2 ,
  \label{eq:dirichlet}
\end{equation}
la contribution de $A_\pi$ s'annulant terme à terme par antisymétrie de $A_\pi$ et symétrie de
$(f_k-f_j)^2$. L'énergie de fluctuation est donc aveugle à la direction.

\emph{(iii) L'irréversibilité tient en un scalaire.} La production d'entropie
\begin{equation}
  \sigma \;=\; \tfrac12\sum_{j,k}\bigl(\pi_j P_{jk}-\pi_k P_{kj}\bigr)
                \log\frac{\pi_j P_{jk}}{\pi_k P_{kj}} \;\geq\; 0
  \label{eq:entropie}
\end{equation}
est nulle si et seulement si la chaîne est réversible, chaque terme étant du signe de
$(x-y)\log(x/y)\geq 0$.

\paragraph{Une précaution de vocabulaire, et elle a été demandée.} Les mots \emph{martingale} et
\emph{variation quadratique} sont employés ici au sens de la théorie des chaînes de Markov, sous
la mesure physique : $\mathcal{E}(f,f)$ est la forme de Dirichlet de la chaîne, c'est-à-dire
la variation quadratique de la martingale associée à la fonction test $f$. Il ne s'agit à aucun
moment d'une mesure risque-neutre, d'un actif répliqué, ni d'un \og{}test de
martingalité\fg{} au sens des générateurs de scénarios économiques. Ce mémoire ne valorise
rien : aucune hypothèse de non-arbitrage n'y est faite ni requise, et aucun changement de
mesure n'y intervient. La confusion est facile, et elle est écartée ici
parce que le même mot désigne, en tarification, un objet qui n'aurait aucun sens. La
martingale, sa filtration et sa mesure sont nommées \matcomp, avec les hypothèses de la
chaîne, la distinction entre ce qui est \emph{testé} et ce qui est \emph{décomposé}, et la
raison pour laquelle aucun test de martingalité n'est possible sur un facteur latent.

\begin{cle}
\textbf{Le placebo, relu, et la portée de cette lecture.} Le non-rejet $z=-0{,}33$ est
compatible avec $\sigma=0$ (équilibre détaillé). Pour une chaîne stationnaire, une statistique
invariante par renversement du temps a la même loi sous $P$ et sous le noyau renversé $\tilde P$ :
elle ne \emph{peut pas} voir un courant qui change de signe sous ce renversement, et c'est ce qui
rend dégénéré le test du \og temps inversé \fg{}. La statistique $T$ du placebo n'est pas de ce
type, puisqu'elle compare les transitions à leur transposée : elle pouvait détecter un courant, et
au pas annuel elle n'en détecte aucun. Deux limites bornent cette lecture. $W$ n'est pas
littéralement un générateur de Markov, et le produit scalaire propre est pondéré par $\pi$, dont le
$S=(W+W^\top)/2$ euclidien est le cas $\pi$ uniforme. Surtout, l'énoncé ne s'étend pas à la loi
complète des ensembles de piliers touchés par un sinistre, qui dépend de la direction
(proposition~\ref{prop:nonid}). La frontière est donc celle des statistiques que les sources
fournissent, non une impossibilité de principe.
\end{cle}

\begin{figure}[htbp]
  \centering
  \figover{Z11_reversibilite_martingale.png}
  \caption{La direction $A$ comme courant irréversible, et
pourquoi une statistique invariante par renversement ne peut pas la voir. (a)~le courant $J_{jk}=\pi_jP_{jk}-\pi_kP_{kj}$
est non nul : le modèle porte bien une direction. (b)~pour trois cents fonctions test tirées au
hasard, l'énergie de fluctuation, c'est-à-dire la variation quadratique de la martingale, est la
même pour la chaîne posée et pour sa symétrisée : les points tombent sur la diagonale, elle ne voit
que $S$. Seule la production d'entropie distingue les deux chaînes, et elle s'annule à la
symétrisation.}
  \label{fig:reversibilite}
\end{figure}
% SOURCES-SCRIPTS: 40

\section{La meilleure expérience naturelle disponible : le choc MOVEit}

Avant de conclure, on tente la stratégie la plus ambitieuse que la donnée autorise : une
expérience quasi naturelle. En mai 2023, l'exploitation de masse de la faille du logiciel de
transfert MOVEit, édité par un tiers, frappe simultanément des centaines
d'organisations. C'est un choc exogène, daté, et typé pilier 4. Si la
contagion dirigée existe, ce choc doit élever, dans les mois suivants et chez les mêmes
organisations, le risque de défaillances relevant d'autres piliers. C'est une prédiction
directionnelle testable, par différence de différences (188 organisations financières touchées,
groupe de contrôle apparié, fenêtres de $180$~jours).

\begin{attention}
\textbf{Le piège, et ce qu'il apprend.} La différence de différences brute est grande et
statistiquement significative ($+1{,}17$ jour-brèche par organisation, IC bootstrap
$[+1{,}07\,;+1{,}27]$). Mais un placebo (le même dispositif calé sur une fausse
date un an plus tôt) produit presque le même chiffre ($+1{,}11$), alors qu'aucun choc n'a
eu lieu. L'effet apparent n'est donc pas causal : il vient de la sélection du groupe de
contrôle, qui revient mécaniquement vers sa moyenne. L'effet net, une fois le
placebo retranché, tombe à $+0{,}06$, indistinguable de zéro. Et quand bien même un signal
subsisterait, la cible resterait non attribuable : la taxonomie de la source est un vecteur
d'attaque (HACK, DISC\dots, $32\,\%$ non déterminés), pas un domaine de contrôle DORA. Même
la meilleure expérience naturelle ne fait donc pas apparaître de propagation dirigée sur
donnée agrégée (figure~\ref{fig:moveit}) : l'échec tient à la donnée, pas à la méthode.
\end{attention}

\begin{figure}[htbp]
\centering
\figover{Z6_event_study_moveit.png}
\caption{Event-study du choc MOVEit. (a)~Fréquence des brèches, groupe touché contre
contrôle, autour du choc. (b)~La différence de différences brute est dévorée par le
placebo : l'effet net est nul. (c)~La cible reste non attribuable, la taxonomie étant un
vecteur d'attaque et non un pilier.}
\label{fig:moveit}
\end{figure}

% SOURCES-SCRIPTS: 62 35

\section{Ce qu'un corpus documentaire corrobore}
\label{sec:postmortems}

Les trois sections précédentes portent sur des bases agrégées. Or la spécification que
$W$ exige, horodatage fin, taxonomie par domaine de contrôle, causes communes consolidées, est
introuvable dans une base agrégée mais existe \emph{à l'intérieur} d'un autre type de source :
le rapport post-incident officiel, où une enquête publique reconstitue la séquence des
défaillances de contrôle. On procède donc à une analyse de contenu structurée de sept
post-mortems (Cyber Safety Review Board, FCA et PRA, GAO, SEC, OCC), dont deux portent sur des
entités financières régulées, en codant chaque rapport en transitions dirigées $k \to j$ telles
que le rapport les \emph{établit}. Le protocole de codage est explicite et réplicable, et
l'unité de codage est le couple (incident, paire de piliers).

Le résultat est net. Sur $17$ transitions codées, l'asymétrie observée vaut $17$ contre un
placebo par permutation de $8{,}4 \pm 2{,}2$, soit $z=+3{,}93$ ($p=0{,}0004$), là où le pas
annuel d'OpRisk donnait $z=-0{,}33$ ; aucune transition ne va à contre-sens de la
tendance dominante de sa paire. Surtout, une structure apparaît :
P1 (gouvernance) n'est jamais une cible ($9$ sorties, $0$ entrée) et P2
(gestion des incidents) est un puits ($0$ sortie, $7$ entrées). C'est exactement la structure
que le classeur qualitatif posait par jugement, retrouvée ici sur donnée documentaire et par un
chemin indépendant. Le compromis est l'inverse de celui d'OpRisk : peu de volume, mais une
direction établie par l'enquête, et c'est la qualité du codage qui donne la puissance.

\begin{attention}
\textbf{Le confondant, qu'aucun de ces chiffres ne lève.} Le placebo teste la \emph{cohérence}
du codage, non l'\emph{exogénéité} de la direction. Or une enquête officielle est \emph{mandatée}
pour établir une cause racine et des responsabilités, et conclut très souvent à une
\og{}défaillance de gouvernance\fg{} : que P1 ressorte comme source pure peut donc refléter la
convention d'écriture des rapports autant que la causalité opérationnelle. S'y ajoutent
un codeur unique, des sources en partie secondaires, et un échantillon raisonné d'incidents
notoires. Ce corpus corrobore donc la structure ; il ne la \emph{calibre} pas. Le
codage à juge unique est à ce titre une limite déclarée de ce mémoire, et non une étape
en attente : son coût est \emph{borné} plutôt que supposé, au \S\ref{sec:biais-narration}, où la
dégradation complète du rôle de $P_1$ se révèle moins chère que l'ignorance directionnelle que
le mémoire assume et publie déjà. Deux dispositifs complètent par ailleurs le dossier, documentés en annexe, et leur exécution
renforcerait l'énoncé sans le sauver : une mesure de fiabilité inter-juges par second
codeur en aveugle, encore \emph{en attente}, et une relecture par praticiens dont les
questions portent précisément sur l'ordre P1 vers P2, à partir de leur expérience de terrain et
non de rapports publics, désormais engagée.

\textbf{Les trois relectures de praticiens reçues vont dans le même sens.} Aucune ne distingue un
ordre de propagation d'une cause commune, ce qui corrobore la \emph{frontière} et non la matrice ;
les deux dernières contestent le sens unique de deux arcs sans déplacer le capital publié,
l'ensemble admissible ne fixant aucune direction (\S\ref{sec:relecture-praticien}).
\end{attention}

% SOURCES-SCRIPTS: 53

\section{Le corpus étendu, et la matrice ordonnée complète}
\label{sec:corpus-etendu}

Deux objections restaient ouvertes. Le signal directionnel pouvait tenir au \emph{choix} des
sept incidents, tous notoires ; et le codage, tel qu'il est exploité ci-dessus, ne compare que
$n(j\to k)$ à $n(k\to j)$ \emph{à l'intérieur} d'une paire, ce qui élimine le dénominateur et
ne dit donc rien de l'amplitude. On lève les deux ensemble.

\paragraph{Le corpus passe à dix rapports, et le signal se renforce.} Trois incidents
s'ajoutent, tous documentés par une enquête publique : CrowdStrike Falcon 2024
(analyse technique de cause racine du 6 août 2024, où le validateur de contenu laisse passer un
fichier dont le nombre de champs ne correspond pas au modèle), Log4Shell 2021--2022
(premier rapport du CSRB, où l'absence d'inventaire des dépendances logicielles
interdit de savoir où la bibliothèque était déployée) et Raphaels Bank 2015 (amende
conjointe FCA/PRA pour manquement répété à la gouvernance de
l'externalisation). Le corpus compte alors $22$ transitions. Le test de direction ne se contente
pas de tenir, il se renforce, et il survit aux deux restrictions les plus dures.

\begin{center}\small
\renewcommand{\arraystretch}{1.25}
\begin{tabular}{@{}l r r r r@{}}
\toprule
\textbf{Périmètre} & \textbf{Rapports} & \textbf{Transitions} & \textbf{Asym.\ / nul} & $z$ \\
\midrule
Corpus initial (chapitre, ci-dessus) & $7$ & $17$ & $17$ / $8{,}4\pm2{,}2$ & $+3{,}93$ \\
\textbf{Corpus étendu} & $10$ & $22$ & $22$ / $9{,}5\pm2{,}4$ & $\mathbf{+5{,}12}$ \\
Sources \emph{primaires} seules & $7$ & $18$ & $18$ / $8{,}2\pm2{,}4$ & $+4{,}08$ \\
Perspective \emph{entité} seule & $6$ & $16$ & $16$ / $8{,}0\pm2{,}2$ & $+3{,}62$ \\
\bottomrule
\end{tabular}
\end{center}

\vspace{2pt}
Les deux dernières lignes comptent autant que la deuxième. Restreindre aux sources
primaires écarte l'hypothèse que le signal viendrait de la couverture de presse ; restreindre à
la perspective entité écarte celle qu'il viendrait du mélange des points de vue, une
défaillance interne d'un fournisseur n'étant pas un risque tiers pour lui-même alors qu'elle
l'est pour ses clients. Le signal survit aux deux.

\paragraph{La matrice ordonnée, enfin estimable.} Chaque incident déclare désormais, en plus de
ses transitions, l'ensemble des piliers défaillants, y compris ceux qui ont cédé sans
rien entraîner. On dispose alors du dénominateur qui manquait, et l'on estime directement
\begin{equation}
  p_{jk} \;=\; \Prob(\text{le pilier } k \text{ cède} \mid \text{le pilier } j \text{ a cédé}),
  \qquad
  p_{jk}\mid\text{données} \;\sim\; \mathrm{Beta}\bigl(1+m_{jk},\;1+N_j-m_{jk}\bigr),
  \label{eq:pjk}
\end{equation}
où $N_j$ compte les incidents où $j$ a cédé et $m_{jk}$ ceux où le rapport établit $j\to k$.
C'est la matrice \emph{ordonnée complète}, donc $S$ et $A$, là où la lecture par paire
ne donnait que le signe de $A$.

\begin{attention}
\textbf{Le résultat est asymétrique, et c'est lui qu'il faut retenir.} Sur les vingt liens
ordonnés, douze ont un crédible à $90\,\%$ entièrement \emph{sous} $0{,}5$, et
\emph{aucun} entièrement au-dessus. Le corpus établit donc très bien \emph{où la contagion ne
va pas}, et pas du tout où elle va fort : même les liens les plus élevés, $P_3\to P_2$
($0{,}77$, crédible $[0{,}48\,;0{,}95]$) et $P_1\to P_4$ ($0{,}68$, $[0{,}40\,;0{,}89]$),
recouvrent encore $0{,}5$. Les douze liens faibles comprennent en revanche les
\emph{quatre} sorties de $P_2$, soit toutes, et trois des quatre entrées de $P_1$ : le puits est
donc établi comme tel, la source pure seulement en partie, l'arc $P_5 \to P_1$ n'étant pas établi
comme faible. Cette asymétrie résiduelle va contre la convention \og{}cause racine =
gouvernance\fg{} plutôt qu'en sa faveur, et elle est rapportée ici pour cette raison.
\end{attention}

\begin{cle}
\textbf{Ce que l'extension confirme, et ce qu'elle ne déplace pas.} La corrélation entre la
matrice documentaire et le classeur d'expert n'est pas significative, ni sur la matrice
complète ($\rho=+0{,}08$) ni sur la seule partie antisymétrique ($\rho=+0{,}31$, $p=0{,}19$).
Le corpus ne reproduit donc pas le classement des coefficients. Ce qu'il reproduit est la structure, et elle seule. Enfin, en
ajoutant des défaillances non propagées au dénominateur, pour tenir compte de ce qu'une enquête
documente mal les contrôles qui cèdent sans conséquence, le \emph{niveau} de $p_{jk}$ s'effondre
(de $0{,}27$ à $0{,}07$ en moyenne) tandis que la part de direction
$\lVert A\rVert/\lVert W\rVert$ ne bouge pas (elle reste à $60\,\%$). Le corpus établit un
ordre, pas une magnitude, et l'extension confirme la frontière à trois niveaux sur un
objet plus riche au lieu de la déplacer.
\end{cle}

% SOURCES-SCRIPTS: 59

\section{Le biais de narration, mesuré au lieu d'être déclaré}
\label{sec:biais-narration}

Des trois limites nommées au \S\ref{sec:postmortems}, la troisième est la plus sérieuse et la
seule qui ne se corrige pas en ajoutant des rapports. Une enquête officielle cherche une
\emph{cause racine}, et \og défaillance de gouvernance\fg{} est une conclusion de convention : le
corpus pourrait surreprésenter $P_1$ comme source non parce que la gouvernance émet, mais parce
que les enquêteurs concluent ainsi. L'enjeu n'est pas mince, puisque onze des vingt-deux
transitions codées partent de $P_1$. Jusqu'ici cette limite était écrite et non chiffrée, ce
qui est la différence entre une réserve et un résultat. On la chiffre.

On la chiffre par une loi nulle exacte, obtenue par convolution et sans tirage, sur laquelle
toutes les décisions se prennent en $p$-valeur exacte plutôt qu'en $z$ ; un jackknife montre
qu'aucun rapport ne porte le signal à lui seul. On cherche ensuite, en borne adverse sur les
$2^{11}$ sous-ensembles d'arêtes partant de $P_1$, l'ampleur qu'il faudrait à la convention
narrative pour renverser la conclusion, en \emph{retirant} ou en \emph{retournant} des arêtes.
La construction et le jackknife sont déportés \matcomp.

\begin{center}\small
\renewcommand{\arraystretch}{1.25}
\begin{tabular}{@{}l l l@{}}
\toprule
\textbf{Hypothèse sur les $11$ arêtes de $P_1$} & \textbf{Reste} & $p$ \textbf{exacte} \\
\midrule
Aucune (corpus tel qu'il est codé) & $22$ transitions & $1{,}5\cdot10^{-5}$ \\
$4$ arêtes \emph{retournées}, pire configuration & $22$ transitions & $\mathbf{0{,}084}$ \\
Les $11$ arêtes \emph{retirées} & $11$ transitions & $0{,}0039$ \\
Les $7$ \emph{rapports} concernés retirés & $3$ transitions & $0{,}50$ \\
\bottomrule
\end{tabular}
\end{center}

\vspace{2pt}
La deuxième ligne donne le point de rupture : il faudrait que la convention narrative ait
\emph{inversé} quatre des onze arêtes, et dans l'agencement le plus défavorable, pour perdre
la significativité. La troisième est plus instructive encore : aucun nombre de retraits
ne suffit. En supprimant la totalité du rôle de la gouvernance, donc la moitié du corpus, la
direction reste significative, portée par les seules chaînes $P_3\to P_2$ et $P_4\to P_2$, les
tests qui ne détectent pas et les tiers qui entraînent la gestion d'incident. Ces deux chaînes
ne doivent rien à la convention \og cause racine = gouvernance\fg{}, puisque $P_1$ n'y figure pas.

La quatrième ligne, qui retire les rapports entiers et non les seules arêtes, jette aussi des
chaînes que le biais ne concerne pas ; c'est une borne de forme, discutée au même endroit.

\begin{cle}
\textbf{Ce que le biais coûte en capital, et ce qu'il ne coûte pas.} Un $z$ n'est pas un
capital. En dégradant l'émission de $P_1$ dans la matrice retenue,
$W_{1k}\to(1-a)\,W_{1k}$, et en \emph{recalculant} la bande d'identification pour chaque degré
(la dégradation modifie $S$, donc l'ensemble admissible lui-même), le SCR d'entité passe
de $169{,}0$ à $127{,}1$~M\euro{} quand $P_1$ devient complètement muette, soit $-24{,}8\,\%$.
Ce déplacement de $41{,}9$~M\euro{} vaut $0{,}81$ fois la largeur de la bande
d'identification déjà publiée, $51{,}7$~M\euro. Accorder au critique la totalité de son
objection coûte donc \emph{moins} cher que l'ignorance directionnelle que le mémoire assume et
publie : le biais de narration n'est pas la première source d'incertitude du chiffre d'entité,
et il ne faut pas le présenter comme telle. Ce qui bouge, en revanche, est le pilotage :
le pilier à remédier en premier passe de $P_1$ à $P_4$ dès $a = 0{,}50$. Le niveau de capital
résiste au biais, la recommandation opérationnelle en dépend, et les deux ne doivent pas être
présentés avec la même assurance.
\end{cle}

Enfin, l'écart entre états de conformité, qui est la thèse du mémoire, survit à tous les degrés
testés : de $+45{,}2\,\%$ à $a=0$ il descend à $+29{,}1\,\%$ à $a=1$ sans s'annuler, la
conformité agissant par le gain $g$ sur \emph{toutes} les arêtes et pas seulement sur celles de
$P_1$. Le second codage en aveugle n'est donc pas le préalable de ce
résultat : il en serait le renforcement.
% SOURCES-SCRIPTS: 54
Ce que le mémoire publie ici est une
\emph{borne adverse}, et le rapport de $0{,}81$ ci-dessus dit pourquoi elle porte : accorder au
critique la totalité de son objection déplace le capital d'entité de moins que l'ignorance
directionnelle déjà assumée. Ce qu'un second codage ajouterait est une mesure
d'\emph{intersubjectivité}, non une correction de niveau, et son absence ne laisse aucun énoncé
du mémoire sans borne.

Le dispositif de second codage est préparé et décrit \matcomp.
% SOURCES-SCRIPTS: 64

\section{Les séquences ordonnées : une voie explorée et refermée}
\label{sec:sequences-ordonnees}

Des marges par paires ne déterminent pas une loi sur les permutations : un triplet ordonné,
du type P1 $\to$ P3 $\to$ P2, porte en général une information qu'aucune paire ne porte. Composer
les arêtes d'un même incident donne $9$ chemins de longueur deux. Le résultat est négatif pour
trois raisons distinctes : la comparaison la plus naturelle a un support vide, P2 n'émettant
jamais et P1 n'étant jamais atteint ; le test de dépendance au prédécesseur n'a aucune puissance,
aucun pilier n'étant à la fois atteint et doté de deux successeurs ; et l'acyclicité du graphe
agrégé s'obtiendrait deux fois sur trois par hasard, donc ne confirme rien. Un acquis subsiste
sur le modèle : l'auto-évitement rend la chaîne des piliers dépendante du prédécesseur, de sorte
que les triplets portent une prédiction testable. Le protocole et le volume requis, $158$ à
$589$ chemins par un même pilier, sont déportés \matcomp.

% SOURCES-SCRIPTS: 75

\section{La formalisation bayésienne, et sa prudence}
\label{sec:bayes}

Borner à la main traite toutes les configurations admissibles à égalité, y compris celles que
les rapports contredisent. Le cadre bayésien fait mieux avec le même matériau, et sans
élicitation. Pour chaque paire $\{lo,hi\}$, on pose $\theta$ la probabilité qu'une chaîne
documentée aille de $lo$ vers $hi$, et on la relie au coefficient antisymétrique par
\[
  A_{lo,hi} \;=\; S_{lo,hi}\,(2\theta - 1),
\]
transformation affine, croissante et bijective de $[0,1]$ sur l'intervalle admissible
$[-S_{lo,hi},+S_{lo,hi}]$ du chapitre~\ref{chap:partielle}, avec $\theta=1/2$ donnant $A=0$,
c'est-à-dire $W$ symétrique. Si le corpus documente $n_{lo}$ chaînes allant de $lo$ vers $hi$ et
$n_{hi}$ dans l'autre sens, la vraisemblance est binomiale et la conjugaison beta-binomiale
donne un posterieur explicite :
\begin{equation}
  \begin{gathered}
  \theta \sim \mathrm{Beta}(\alpha_0,\beta_0),
  \qquad
  n_{lo}\mid\theta \sim \mathcal{B}\mathrm{in}(n_{lo}+n_{hi},\,\theta)
  \\
  \Longrightarrow\qquad
  \theta \mid \text{données} \;\sim\; \mathrm{Beta}\bigl(\alpha_0+n_{lo},\;\beta_0+n_{hi}\bigr).
  \end{gathered}
  \label{eq:posterieur-beta}
\end{equation}
Aucune simulation par chaînes de Markov n'est nécessaire, et aucun prior d'expert n'intervient.
Le cas $n_{lo}=n_{hi}=0$, celui des paires que le corpus ne documente pas, rend
\eqref{eq:posterieur-beta} égal au prior : sous $\mathrm{Beta}(1,1)$, $\theta$ reste uniforme sur
$[0,1]$ et $A$ balaie donc \emph{tout} l'intervalle admissible.

\begin{cle}
\textbf{La propriété qui fait l'intérêt du procédé.} Le posterieur reproduit de
lui-même la structure d'identification partielle : les quatre paires que le corpus ne documente
pas gardent un $\theta$ uniforme, donc un coefficient qui couvre \emph{tout} l'intervalle
admissible, sans qu'on ait à poser cette ignorance ; les paires documentées, elles, se
concentrent. Le bayésien n'ajoute donc pas une hypothèse, il formalise celle qui était déjà là.
Le capital devient une loi plutôt qu'un encadrement, et son intervalle de crédibilité à
$90\,\%$ se révèle $81\,\%$ plus resserré que les bornes de pire cas (les niveaux sont établis
en partie~\ref{part:resultats}). Les deux objets répondent à deux questions distinctes, et le
mémoire donne les deux : les bornes sont \emph{opposables} car sans prior, le crédible est
\emph{informatif}.
\end{cle}

Deux vérifications vont en sens opposés. Sous un prior sceptique, la médiane ne se déplace
presque pas, donc la conclusion vient des comptes documentés. Mais par paire, seules deux des six
paires documentées ont un intervalle de crédibilité à $90\,\%$ qui exclut $\theta=1/2$ : c'est
cette lecture, la plus conservatrice, qui est retenue. Le détail et la figure sont déportés
\matcomp.
% SOURCES-SCRIPTS: 56

\section{La frontière, en trois niveaux}

On peut maintenant énoncer précisément ce que la donnée autorise à revendiquer, et c'est un
énoncé à trois étages plutôt qu'un verdict unique.

\begin{center}\small
\renewcommand{\arraystretch}{1.3}
\begin{tabular}{@{}p{3.2cm} p{7.4cm} p{4.2cm}@{}}
\toprule
\textbf{Objet} & \textbf{Ce que la donnée établit} & \textbf{Posture retenue} \\
\midrule
\textbf{Structure} \newline (P1 amont, P2 aval) &
$17$ transitions sur le corpus initial et $22$ sur le corpus étendu, aucune à contre-sens,
$z=+3{,}9$ puis $+5{,}1$ ; convergence avec le jugement d'expert par un chemin indépendant &
\textbf{revendiquée}, avec son confondant déclaré \\
\textbf{Direction} \newline paire par paire &
$2$ paires sur $6$ seulement ont un crédible concluant &
\textbf{non revendiquée} ; les bornes restent \\
\textbf{Amplitude} \newline des coefficients &
aucune information ; même les dix sens connus, une bande de $656$~M\euro{} subsisterait &
\textbf{non identifiable} par cette voie \\
\bottomrule
\end{tabular}
\end{center}

\vspace{4pt}
La conséquence de méthode est nette : on revendique un ordre partiel entre piliers, non
une matrice. C'est un objet plus faible, donc bien plus difficile à contester, et il \emph{suffit}
à porter les conclusions du mémoire, à savoir le classement des piliers et l'ordre de
remédiation. Ce que la donnée fixe par ailleurs demeure ce qu'établissent les chapitres
empiriques : la fréquence marginale d'entrée, l'accumulation par cause commune, et l'indice de
queue de sévérité avec son hétérogénéité par catégorie. Ce qu'elle ne fixe pas relève d'un autre
régime de connaissance et se lit en bande, jamais en point.

La spécification du registre qui identifierait $W$, et la valeur en capital de cette
information, sont traitées au chapitre~\ref{ch:livrable}, en partie~\ref{part:limites} : elles
relèvent de la portée réglementaire du travail plus que de sa démonstration.

\begin{cle}
\textbf{Ce que ce chapitre apporte au mémoire.} Démontrer, chiffres à l'appui, qu'une dépendance
que tout le monde \emph{postule} n'est pas \emph{mesurable} sur les données agrégées est rare ;
en donner la raison théorique l'est davantage ; aller ensuite chercher, dans une source d'un
autre genre, de quoi en corroborer la structure sans sur-vendre le résultat est ce qui distingue
une limite \emph{assumée} d'une limite \emph{subie}. Il en sort quatre acquis : une
\emph{frontière d'identification} à trois niveaux, une discipline sur l'usage de $W$
(ordre partiel revendiqué, coefficients bornés, jamais présentés comme calibrés), une
\emph{formalisation} bayésienne qui reproduit cette frontière au lieu de la postuler, et une
spécification de reporting hiérarchisée qui donne au travail une portée réglementaire
au-delà du modèle.
\end{cle}

% =====================================================================

% SOURCES-SCRIPTS: 55 59 64

